\chapter[Vermeintliche Rettung]{Vermeintliche Rettung}
\quellseite{11}
Vor einigen Jahren gerieth ein Boot, in dem zwei Menschen saßen, oberhalb des Niagarafalles in den Strudel. In wenigen Augenblicken war es umgestürzt, und die armen Menschen rangen mit den Fluthen. Unaufhaltsam trieben sie dem Abgrund zu. Am Ufer standen entsetzte Menschen, die die schreckliche Lage der Unglücklichen beobachteten, ohne zu wissen, wie sie ihnen helfen sollten. Da eilte ein Mann mit einem Seil herbei und warf es ihnen zu. Der eine von ihnen ergriff es sofort. Er wurde gerettet. In demselben Augenblick kam ein großer Holzblock angeschwommen. An ihn klammerte sich der zweite und versäumte darüber das Ergreifen des Rettungsseils.

Als er seinen Mißgriff einsah, war es zu spät. Er stürzte hülflos sammt dem Block in die Tiefe.

Jesu ist das Rettungsseil. Nur wer ihn ergriffen hat, ist gerettet. Wer irgend eine andere, scheinbar noch so sichere Zuflucht hat, ist rettungslos verloren.
